\documentclass[spanish,12pt,a4paper]{article}
\usepackage[T1]{fontenc}
\usepackage{babel}
\usepackage{fullpage}
\usepackage{termcal}
\usepackage{xcolor}
\usepackage{here}
\usepackage{hyperref}
\usepackage{booktabs}
\usepackage{here}
\usepackage{array}
\usepackage{multirow, makecell}

\title{Data Structures}
\date{Course 2025-2026}


\renewcommand{\calprintdate}{%
    \framebox{\monthname\ \ordinaldate}%
}

% Few useful commands (our classes always meet either on Monday and Wednesday 
% or on Tuesday and Thursday)

\newcommand{\TFClass}{%
\calday[Monday]{\classday} % Monday 
\skipday % Tuesday 
\skipday  % Wednesday
\calday[Thursday]{\classday} % Thursday
\skipday % Friday 
\skipday\skipday % weekend (no class)
}

%\newcommand{\TRClass}{%
%\skipday % Monday (no class)
%\calday[Tuesday]{\classday} % Tuesday
%\skipday % Wednesday (no class)
%\calday[Thursday]{\classday} % Thursday
%\skipday % Friday 
%\skipday\skipday % weekend (no class)
%}

\newcommand{\Holiday}[2]{%
\options{#1}{\noclassday}
\caltext{#1}{#2}
}

\begin{document}
\maketitle
\subsection*{Theory and Lab lectures:}
\begin{description}
\item[Monday:] 12:00-14:00 $\rightarrow$ Theory $\rightarrow$ Classroom: A-020
\item[Thursday:] 12:00-14:00 $\rightarrow$ Lab $\rightarrow$ Classroom: A-331
\end{description}

\subsection*{Lecturer:}
\noindent Ugaitz Amozarrain\\
Department of Statistics, Computer Science and Mathematics\\
``Las Encinas'' building, basement, office 9002\\
\href{mailto:ugaitz.amozarrain@unavarra.es}{ugaitz.amozarrain@unavarra.es} 

\subsection*{Office hours:}
\begin{itemize}
\item Monday: 16:00-18:00
\item Wednesday: 11:00-13:00
\end{itemize}

By default, otherwise stated, all office hours will be in my office, though it will be helpful if we agree on the time slot. 
You can contact me either through the internal email of MiAulario, or at \href{mailto:ugaitz.amozarrain@unavarra.es}{ugaitz.amozarrain@unavarra.es}.
Thus, we coordinate the meetings and avoid crowds in the department corridor.

\newpage

\subsection*{Tentative Schedule:}
\begin{center}
% CAREFUL the date is in american month/day/year
\begin{calendar}{1/26/2026}{18} % Semester starts on 1/26/2026 and last for 18 weeks, including finals week
                    %TODO adjust weeks
\setlength{\calboxdepth}{.3in}
\TFClass
% schedule
\caltexton{1}{Presentation}
\caltextnext{}
\caltextnext{Review Python basics / Tips and tricks}
\caltextnext{Sparse Matrices}
\caltextnext{Object-Oriented Python}
\caltextnext{Working with Files}
\caltextnext{Lists, Dictionaries, Stacks}
\caltextnext{Postfix Calculator}
\caltextnext{Queues and Dequeues}
\caltextnext{Palindromes}
\caltextnext{Buffers}
\caltextnext{Balanced Parentheses}
\caltextnext{Linked Lists I}
\caltextnext{Conway's Game of Life}
\caltextnext{Linked Lists II}
\caltextnext{Linked Lists III}
\caltextnext{Working with Linked Lists}
\caltextnext{Recursion and Computational Complexity}
\caltextnext{Binary Trees}
\caltextnext{Recursion Exercises}
\caltextnext{Constructing Binary Trees}
\caltextnext{Generalizing Trees}
\caltextnext{Huffman Coding I}
\caltextnext{Graphs}
\caltextnext{Huffman Coding II}
\caltextnext{Review}
\caltextnext{Huffman Coding III}
% ... and so on

% Holidays
\Holiday{3/19/2026}{\textcolor{red}{\textbf{{\large Father's Day}}}}
\Holiday{4/2/2026}{\textcolor{red}{\textbf{{\large Easter}}}}
\Holiday{4/6/2026}{\textcolor{red}{\textbf{{\large Easter}}}}
\Holiday{4/9/2026}{\textcolor{red}{\textbf{{\large Easter}}}}
\Holiday{4/20/2026}{\textcolor{red}{\textbf{{Day of the University}}}}
% ... and so on
% TODO make it show finals week
\options{5/18/2026}{\noclassday} % finals week
\options{5/21/2026}{\noclassday} % finals week
\options{5/25/2026}{\noclassday} % finals week
\options{5/28/2026}{\noclassday} % finals week
\end{calendar}
\end{center}


\subsection*{Exams:}
% TODO check dates
\begin{description}
\item[25/05/2026:] 16:00 $\rightarrow$ Final exam: Everything
\item[17/06/2026:] 16:00 $\rightarrow$ Retake exam: Everything
\end{description}

\subsection*{Evaluation criteria:}

\begin{table}[H]
\def\arraystretch{1.5}
\begin{tabular}{|p{0.6\linewidth}|c|c|}
\hline
\textbf{Evaluation method} & \textbf{Weigth} (\%) & \textbf{Retake policy} \\
\hline
\hline
Theoretical/Practical exams (Individually performed). Every student must get at least 50\% on the exams to average with the rest of the gradable aspects. & 50\% & YES \\
\hline
Deliverables & 50\% & YES \\
\hline
\end{tabular}
\end{table}

\newpage


\begin{itemize}
\item Exam (5 points):
\begin{itemize}
\item It is mandatory to pass the exam in order to pass the course.
\item The re-take exam is worth 100\% of the grade. 
\end{itemize}
\item Deliverables (5 points):
\begin{itemize}
\item All deliverables will be checked to confirm correct delivery.
\item Each week you will have to deliver the code for the corresponding lab project.
\\
\vspace{-0.5cm}
\begin{table}[h!btp]
\renewcommand{\arraystretch}{1.2}
\centering
\begin{tabular}{m{0.35\textwidth}@{\hskip 0.05\textwidth}m{0.3\textwidth}m{0.2\textwidth}}
\toprule
\textbf{Lab Project} & \textbf{Delivery/Correction} & \textbf{Max points} \\
\midrule
Linked Lists & Correction & 2 \\
Huffman Coding & Correction & 2 \\
\midrule
Sparse Matrices & Delivery & \multirowcell{8}{{\Huge $\frac{n}{8}$}, where n is \\the number of \\delivered items} \\
Working with Files & Delivery &  \\
Postfix Calculator & Delivery  &  \\
Palindromes & Delivery &  \\
Balanced Parentheses & Delivery &  \\
Conway's game of Life & Delivery & \\
Recursion Exercises & Delivery &  \\
Constructing Binary Trees & Delivery &  \\
\bottomrule
\end{tabular}
\end{table}
\end{itemize}
\end{itemize}




\end{document}
